\documentclass[a4paper,12pt,twocolumn]{article}

\usepackage[landscape,top=1cm,bottom=1.4cm,left=1cm,right=1cm,columnsep=1cm]{geometry}
\usepackage[fleqn]{amsmath}
\usepackage{amssymb}
\setlength{\mathindent}{1.5em}
\usepackage{fontspec}
\usepackage{enumitem}
\usepackage{framed}
\usepackage{needspace}
\usepackage[hidelinks]{hyperref}

% Set fonts
\setmainfont{Times New Roman}
\setsansfont{Arial}

\pagestyle{plain}
\setlength{\columnseprule}{0.2pt}
\newcommand{\examdate}{2026/03/26}
\newlength{\problemnumberwidth}
\settowidth{\problemnumberwidth}{\textbf{94.}\ }
\newlength{\subproblemnumberwidth}
\settowidth{\subproblemnumberwidth}{\textbf{57(a).}\ }
\newenvironment{solutionbox}{\begin{framed}}{\end{framed}}

\newcommand{\sectiontitle}[2]{%
  \hypertarget{#1}{}%
  \par\vspace{0.8em}%
  \noindent{\Large\bfseries #2}\par
  \vspace{0.3em}%
  \hrule
  \vspace{0.8em}%
}

\newcommand{\tocproblem}[2]{%
  \hyperref[#1]{#2} & \hyperref[#1]{\pageref*{#1}} \\
}

\newcommand{\worksheettoc}{%
  \noindent{\Large\bfseries Contents}\par
  \vspace{0.3em}%
  \hrule
  \vspace{0.9em}%
  \begin{tabular}{@{}p{0.72\columnwidth}r@{}}
    \textbf{Problem} & \textbf{Page} \\
    \tocproblem{prob:10-1-16}{10-1 \quad 16}
    \tocproblem{prob:10-1-19}{10-1 \quad 19}
    \tocproblem{prob:10-1-22}{10-1 \quad 22}
    \tocproblem{prob:10-1-57a}{10-1 \quad 57(a)}
    \tocproblem{prob:10-1-57b}{10-1 \quad 57(b)}
    \tocproblem{prob:10-4-62}{10-4 \quad 62}
    \tocproblem{prob:11-3-34}{11-3 \quad 34}
    \tocproblem{prob:11-10-1}{11-10 \quad 1}
    \tocproblem{prob:11-10-2}{11-10 \quad 2}
    \tocproblem{prob:11-10-93}{11-10 \quad 93}
    \tocproblem{prob:11-10-94}{11-10 \quad 94}
  \end{tabular}\par
}

\newcommand{\problem}[3]{%
  \Needspace{10\baselineskip}%
  \phantomsection\label{#1}%
  \par\vspace{0.8em}%
  \noindent
  \hangindent=\problemnumberwidth
  \hangafter=1
  \makebox[\problemnumberwidth][l]{\textbf{#2.}}#3\par
}

\newcommand{\subproblem}[3]{%
  \Needspace{10\baselineskip}%
  \phantomsection\label{#1}%
  \par\vspace{0.8em}%
  \noindent
  \hangindent=\subproblemnumberwidth
  \hangafter=1
  \makebox[\subproblemnumberwidth][l]{\textbf{#2.}}#3\par
}

\newlist{subparts}{enumerate}{1}
\setlist[subparts]{%
  label=(\alph*),
  labelindent=\problemnumberwidth,
  leftmargin=*,
  labelsep=0.5em,
  align=left,
  itemsep=0.2em,
  topsep=0.3em,
  parsep=0pt,
  partopsep=0pt
}

\newlist{steps}{enumerate}{1}
\setlist[steps]{%
  label=\textbf{Step \arabic*.},
  leftmargin=*,
  itemsep=0.35em,
  topsep=0.35em,
  align=left
}

\begin{document}

\twocolumn[
{\centering
\Large\bfseries Calculus Problems Solutions\par
\vspace{0.5em}
\noindent\hfill\normalsize\normalfont \examdate\par
\vspace{0.8em}
}
]

\worksheettoc
\vfill\null
\newpage

\sectiontitle{sec:10-1}{10-1}

\problem{prob:10-1-16}{16}{%
\begin{subparts}
\item Eliminate the parameter to find a Cartesian equation of the curve.
\item Sketch the curve and indicate with an arrow the direction in which the curve is traced as the parameter increases.
\end{subparts}
\[
x=\csc t,\ y=\cot t,\ 0<t<\pi
\]}

\begin{solutionbox}
\textbf{Method 1: Use the identity \(\csc^2 t-\cot^2 t=1\)}

\begin{steps}
\item The two coordinates are already written in terms of \(\csc t\) and \(\cot t\), so the identity
\[
\csc^2 t-\cot^2 t=1
\]
is the most direct elimination tool. This step is appropriate because it removes the parameter without solving for \(t\).

Substituting \(x=\csc t\) and \(y=\cot t\), we get
\[
x^2-y^2=1.
\]

\item The equation \(x^2-y^2=1\) is a hyperbola, but the parameter restriction tells us which branch actually appears. Since \(0<t<\pi\), we have \(\sin t>0\), so
\[
x=\csc t>0.
\]
Therefore only the \textbf{right-hand branch} is traced.

\item To determine the direction, inspect easy values of \(t\):
\[
t\to 0^+ \Longrightarrow \sin t\to 0^+, \ \cot t\to +\infty,
\]
so \((x,y)\to(\infty,\infty)\).
Also,
\[
t=\frac{\pi}{2}\Longrightarrow (x,y)=(1,0),
\]
and
\[
t\to \pi^- \Longrightarrow (x,y)\to(\infty,-\infty).
\]
These checkpoints show that the point moves downward along the right branch.
\end{steps}
\end{solutionbox}

\begin{solutionbox}
\textbf{Method 2: Solve for \(\sin t\) and \(\cos t\) in terms of \(x\) and \(y\)}

\begin{steps}
\item Because \(x=\csc t\), taking the reciprocal gives
\[
\sin t=\frac{1}{x}.
\]
This is a natural first move because the reciprocal of \(\csc t\) is simple, and it prepares the standard identity \(\sin^2 t+\cos^2 t=1\).

\item Next find \(\cos t\) from the ratio of the coordinates:
\[
\frac{y}{x}=\frac{\cot t}{\csc t}=\cos t.
\]
This step is appropriate because it expresses the second basic trigonometric function entirely in terms of \(x\) and \(y\).

\item Now apply
\[
\sin^2 t+\cos^2 t=1:
\]
\[
\left(\frac{1}{x}\right)^2+\left(\frac{y}{x}\right)^2=1.
\]
Multiplying by \(x^2\) gives
\[
1+y^2=x^2,
\]
so
\[
x^2-y^2=1.
\]

\item The restriction \(0<t<\pi\) again implies \(x>0\), so the traced part is the right branch only. Since \(\cot t\) decreases from \(+\infty\) to \(-\infty\) on \((0,\pi)\), the motion goes from the upper-right portion through \((1,0)\) to the lower-right portion.
\end{steps}
\end{solutionbox}

\begin{solutionbox}
\textbf{Hand-in Version}

\[
\boxed{x^2-y^2=1}
\]

Because \(0<t<\pi\), we have \(\sin t>0\), so \(x=\csc t>0\). Therefore the curve is the \textbf{right branch} of the hyperbola \(x^2-y^2=1\).

As \(t\) increases,
\[
(\infty,\infty)\to(1,0)\to(\infty,-\infty),
\]
so the direction is from upper-right to lower-right.
\end{solutionbox}

\problem{prob:10-1-19}{19}{%
\begin{subparts}
\item Eliminate the parameter to find a Cartesian equation of the curve.
\item Sketch the curve and indicate with an arrow the direction in which the curve is traced as the parameter increases.
\end{subparts}
\[
x=\ln t,\ y=\sqrt{t},\ t\geq 1
\]}

\begin{solutionbox}
\textbf{Method 1: Solve for \(t\) from the logarithm}

\begin{steps}
\item The equation \(x=\ln t\) is easy to invert because the logarithm is one-to-one on \(t>0\). That makes it the cleanest place to start.

Exponentiating gives
\[
t=e^x.
\]

\item Substitute this into \(y=\sqrt{t}\):
\[
y=\sqrt{e^x}=e^{x/2}.
\]
Equivalently,
\[
y^2=e^x.
\]
This removes the parameter and produces the Cartesian relation directly.

\item The condition \(t\ge 1\) becomes
\[
x=\ln t\ge 0
\qquad\text{and}\qquad
y=\sqrt{t}\ge 1.
\]
So the graph begins at
\[
t=1 \Longrightarrow (x,y)=(0,1),
\]
and only the part with \(x\ge 0\) is used.

\item As \(t\) increases, both \(\ln t\) and \(\sqrt t\) increase, so the point moves steadily up and to the right. This gives the correct tracing direction.
\end{steps}
\end{solutionbox}

\begin{solutionbox}
\textbf{Method 2: Solve for \(t\) from the square root}

\begin{steps}
\item From
\[
y=\sqrt t,
\]
we get
\[
t=y^2.
\]
This is appropriate because it turns the parameter into a very simple algebraic expression.

\item Substitute \(t=y^2\) into \(x=\ln t\):
\[
x=\ln(y^2)=2\ln y.
\]
Solving for \(y\) gives
\[
y=e^{x/2}.
\]
So this method reaches the same Cartesian equation from the opposite direction.

\item Since \(t\ge 1\), we have \(y\ge 1\). Then
\[
x=2\ln y\ge 0,
\]
so the curve again starts at \((0,1)\) and lies only in the region \(x\ge 0\), \(y\ge 1\).

\item The parameter direction is also easy to read here: larger \(t\) means larger \(y\), and then \(x=\ln t\) increases as well. Hence the motion is upward and to the right.
\end{steps}
\end{solutionbox}

\begin{solutionbox}
\textbf{Hand-in Version}

\[
\boxed{y=e^{x/2}}
\qquad\text{or equivalently}\qquad
\boxed{y^2=e^x}
\]

Because \(t\ge 1\), the curve starts at
\[
(x,y)=(0,1)
\]
and only the part with \(x\ge 0\), \(y\ge 1\) is traced.

As \(t\) increases, both coordinates increase, so the point moves up and to the right.
\end{solutionbox}

\problem{prob:10-1-22}{22}{%
\begin{subparts}
\item Eliminate the parameter to find a Cartesian equation of the curve.
\item Sketch the curve and indicate with an arrow the direction in which the curve is traced as the parameter increases.
\end{subparts}
\[
x=\sinh t,\ y=\cosh t
\]}

\begin{solutionbox}
\textbf{Method 1: Use the hyperbolic identity}

\begin{steps}
\item The coordinates are \(\sinh t\) and \(\cosh t\), so the matching identity
\[
\cosh^2 t-\sinh^2 t=1
\]
is the direct elimination formula. It is the natural choice because it contains exactly the two functions in the problem.

\item Substitute \(x=\sinh t\) and \(y=\cosh t\):
\[
y^2-x^2=1.
\]
So the Cartesian equation is a hyperbola.

\item To identify the traced branch, use the range of \(\cosh t\). Since
\[
\cosh t\ge 1 \quad \text{for all real } t,
\]
we must have
\[
y\ge 1.
\]
Therefore the curve is the \textbf{upper branch} of the hyperbola.

\item For direction, test key values:
\[
t\to-\infty \Longrightarrow (x,y)\to(-\infty,\infty),
\]
\[
t=0 \Longrightarrow (x,y)=(0,1),
\]
\[
t\to\infty \Longrightarrow (x,y)\to(\infty,\infty).
\]
So as \(t\) increases, the point moves from left to right along the upper branch.
\end{steps}
\end{solutionbox}

\begin{solutionbox}
\textbf{Method 2: Use the exponential definitions}

\begin{steps}
\item Write
\[
\sinh t=\frac{e^t-e^{-t}}{2},
\qquad
\cosh t=\frac{e^t+e^{-t}}{2}.
\]
This is appropriate because adding and subtracting these expressions isolates \(e^t\) and \(e^{-t}\).

\item Compute
\[
y+x=\cosh t+\sinh t=e^t,
\qquad
y-x=\cosh t-\sinh t=e^{-t}.
\]
Now the parameter appears only in reciprocal exponentials, which is ideal for elimination.

\item Multiply the two equations:
\[
(y+x)(y-x)=e^t e^{-t}=1.
\]
Thus
\[
y^2-x^2=1.
\]

\item Since \(y=\cosh t\ge 1\), only the upper branch is traced. Also, \(e^t\) increases and \(e^{-t}\) decreases as \(t\) increases, so the point slides from the left side of the branch through \((0,1)\) and out to the right side.
\end{steps}
\end{solutionbox}

\begin{solutionbox}
\textbf{Hand-in Version}

\[
\boxed{y^2-x^2=1}
\]

Because \(y=\cosh t\ge 1\), the curve is the \textbf{upper branch} of the hyperbola \(y^2-x^2=1\).

As \(t\) increases, the point moves from the upper-left through \((0,1)\) to the upper-right.
\end{solutionbox}

\problem{prob:10-1-57}{57}{Find the point at which the parametric curve intersects itself and the corresponding values of \(t\).}

\subproblem{prob:10-1-57a}{57(a)}{\(x=1-t^2,\ y=t-t^3\)}

\begin{solutionbox}
\textbf{Method 1: Use the even-odd symmetry}

\begin{steps}
\item Observe that
\[
x(t)=1-t^2
\]
is an even function, while
\[
y(t)=t-t^3
\]
is an odd function. This is useful because self-intersections often occur when the same \(x\)-value is produced by \(t\) and \(-t\).

\item Since
\[
x(-t)=x(t),
\]
the remaining requirement for \(t\) and \(-t\) to describe the same point is
\[
y(-t)=y(t).
\]
But \(y\) is odd, so
\[
y(-t)=-y(t).
\]
Therefore equality can happen only when
\[
y(t)=0.
\]

\item Solve
\[
t-t^3=t(1-t^2)=0.
\]
So
\[
t=0,\ \pm 1.
\]
The value \(t=0\) does not create a self-intersection because it gives only one parameter value. The distinct pair is
\[
t=1 \quad \text{and} \quad t=-1.
\]

\item Evaluate the point:
\[
x=1-(\pm1)^2=0,
\qquad
y=(\pm1)-(\pm1)^3=0.
\]
Hence the curve intersects itself at
\[
(0,0).
\]
\end{steps}
\end{solutionbox}

\begin{solutionbox}
\textbf{Method 2: Solve the two-parameter system directly}

\begin{steps}
\item Let \(t_1\neq t_2\) produce the same point. Then both coordinates must match:
\[
1-t_1^2=1-t_2^2,
\qquad
t_1-t_1^3=t_2-t_2^3.
\]
This is the formal definition of a self-intersection, so it is the systematic place to start.

\item From the \(x\)-equation,
\[
t_1^2=t_2^2,
\]
so
\[
t_2=\pm t_1.
\]
Because \(t_1\neq t_2\), the case \(t_2=t_1\) is impossible. Therefore
\[
t_2=-t_1.
\]

\item Substitute this into the \(y\)-equation:
\[
t_1-t_1^3=-t_1+t_1^3.
\]
Thus
\[
2t_1-2t_1^3=0
\quad\Longrightarrow\quad
2t_1(1-t_1^2)=0.
\]
So
\[
t_1=0,\ \pm1.
\]
Again \(t_1=0\) would force \(t_2=0\), which is not a self-intersection. Hence the valid pair is
\[
t_1=1,\qquad t_2=-1.
\]

\item Substituting either value gives
\[
(x,y)=(0,0).
\]
\end{steps}
\end{solutionbox}

\begin{solutionbox}
\textbf{Hand-in Version}

The curve intersects itself at
\[
\boxed{(0,0)}.
\]

The corresponding parameter values are
\[
\boxed{t=1 \text{ and } t=-1.}
\]

To justify that this is truly a self-intersection, notice that the two
\emph{distinct} parameter values \(t=1\) and \(t=-1\) give the same point:
\[
x(1)=1-1=0,\qquad y(1)=1-1=0,
\]
\[
x(-1)=1-1=0,\qquad y(-1)=-1-(-1)^3=0.
\]
So the curve passes through \((0,0)\) twice.

Moreover,
\[
\frac{dy}{dx}=\frac{dy/dt}{dx/dt}
=\frac{1-3t^2}{-2t}.
\]
At \(t=1\),
\[
\frac{dy}{dx}=\frac{-2}{-2}=1,
\]
while at \(t=-1\),
\[
\frac{dy}{dx}=\frac{-2}{2}=-1.
\]
The tangent slopes are different, so the two branches cross each other at the
origin. Therefore \((0,0)\) is genuinely an intersection point of the graph.
\end{solutionbox}

\subproblem{prob:10-1-57b}{57(b)}{\(x=2t-t^3,\ y=t-t^2\)}

\begin{solutionbox}
\textbf{Method 1: Start with the simpler equation \(y(t_1)=y(t_2)\)}

\begin{steps}
\item Let \(t_1\neq t_2\) give the same point. Then
\[
t_1-t_1^2=t_2-t_2^2.
\]
This is the better equation to start from because it is only quadratic, while the \(x\)-equation is cubic.

\item Rearranging gives
\[
(t_1-t_2)-(t_1^2-t_2^2)=0.
\]
Factor:
\[
(t_1-t_2)\bigl[1-(t_1+t_2)\bigr]=0.
\]
Since \(t_1\neq t_2\), we must have
\[
t_1+t_2=1.
\]

\item Now use the equality of the \(x\)-coordinates:
\[
2t_1-t_1^3=2t_2-t_2^3.
\]
Move everything to one side and factor:
\[
2(t_1-t_2)-(t_1^3-t_2^3)=0,
\]
\[
(t_1-t_2)\bigl[2-(t_1^2+t_1t_2+t_2^2)\bigr]=0.
\]
Again \(t_1\neq t_2\), so
\[
t_1^2+t_1t_2+t_2^2=2.
\]

\item Let
\[
s=t_1+t_2,\qquad p=t_1t_2.
\]
Because \(s=1\),
\[
t_1^2+t_1t_2+t_2^2=s^2-p=1-p.
\]
So
\[
1-p=2
\quad\Longrightarrow\quad
p=-1.
\]
Therefore \(t_1,t_2\) are the roots of
\[
u^2-su+p=0
\quad\Longrightarrow\quad
u^2-u-1=0.
\]
Hence
\[
t=\frac{1\pm\sqrt5}{2}.
\]

\item To find the common point, use the relation \(t^2=t+1\). Then
\[
y=t-t^2=t-(t+1)=-1.
\]
Also,
\[
t^3=t(t^2)=t(t+1)=t^2+t=(t+1)+t=2t+1,
\]
so
\[
x=2t-t^3=2t-(2t+1)=-1.
\]
Thus the self-intersection point is
\[
(-1,-1).
\]
\end{steps}
\end{solutionbox}

\begin{solutionbox}
\textbf{Method 2: Use symmetric sums from the beginning}

\begin{steps}
\item Suppose \(t_1\neq t_2\) produce the same point, and set
\[
s=t_1+t_2,\qquad p=t_1t_2.
\]
This is appropriate because factorization naturally turns the equal-coordinate equations into statements about \(s\) and \(p\).

\item From
\[
t_1-t_1^2=t_2-t_2^2
\]
we get
\[
(t_1-t_2)\bigl[1-(t_1+t_2)\bigr]=0.
\]
Since \(t_1\neq t_2\), it follows that
\[
s=1.
\]

\item From
\[
2t_1-t_1^3=2t_2-t_2^3
\]
we obtain
\[
(t_1-t_2)\bigl[2-(t_1^2+t_1t_2+t_2^2)\bigr]=0.
\]
So
\[
t_1^2+t_1t_2+t_2^2=2.
\]
In terms of \(s\) and \(p\),
\[
t_1^2+t_1t_2+t_2^2=s^2-p.
\]
Because \(s=1\), this gives
\[
1-p=2
\quad\Longrightarrow\quad
p=-1.
\]

\item Therefore \(t_1,t_2\) are the roots of
\[
u^2-u-1=0,
\]
so
\[
t=\frac{1\pm\sqrt5}{2}.
\]

\item Each root satisfies \(t^2=t+1\), so
\[
y=t-t^2=-1,
\qquad
x=2t-t^3=-1.
\]
Hence the common point is
\[
(-1,-1).
\]
\end{steps}
\end{solutionbox}

\begin{solutionbox}
\textbf{Hand-in Version}

The curve intersects itself at
\[
\boxed{(-1,-1)}.
\]

The corresponding parameter values are
\[
\boxed{t=\frac{1+\sqrt5}{2}\quad\text{and}\quad t=\frac{1-\sqrt5}{2}.}
\]

To justify that the graph really intersects itself, let
\[
t_1=\frac{1+\sqrt5}{2},\qquad t_2=\frac{1-\sqrt5}{2}.
\]
These are distinct real numbers, and both satisfy
\[
t^2=t+1.
\]
Hence for either one,
\[
y=t-t^2=t-(t+1)=-1,
\]
and
\[
t^3=t(t^2)=t(t+1)=2t+1,
\]
so
\[
x=2t-t^3=2t-(2t+1)=-1.
\]
Thus
\[
\bigl(x(t_1),y(t_1)\bigr)=\bigl(x(t_2),y(t_2)\bigr)=(-1,-1),
\]
which proves that the same point is reached by two different parameter values.

Also,
\[
\frac{dy}{dx}=\frac{dy/dt}{dx/dt}
=\frac{1-2t}{2-3t^2}.
\]
Using \(t^2=t+1\), this becomes
\[
\frac{dy}{dx}=\frac{1-2t}{-1-3t}.
\]
So at \(t=t_1\),
\[
\frac{dy}{dx}=\frac{3-\sqrt5}{2},
\]
and at \(t=t_2\),
\[
\frac{dy}{dx}=\frac{3+\sqrt5}{2}.
\]
These slopes are different, so the two branches cross at \((-1,-1)\) rather
than merely touching there.
\end{solutionbox}

\sectiontitle{sec:4}{10-4}

\problem{prob:10-4-62}{62}{Graph the polar curve. Choose a parameter interval that produces the entire curve.
\[
r=\lvert\tan\theta\rvert^{\lvert\cot\theta\rvert}.
\]
\noindent(valentine curve)}

\begin{solutionbox}
\textbf{Method 1: Rewrite the radius as \(u^{1/u}\)}

\begin{steps}
\item Let
\[
u=\lvert\tan\theta\rvert.
\]
Then \(u>0\) whenever the formula is defined, and
\[
\lvert\cot\theta\rvert=\frac{1}{u}.
\]
So the polar equation becomes
\[
r=u^{1/u}.
\]
This substitution is appropriate because it turns the unusual polar formula into a one-variable function whose size and extrema are easier to study.

\item To understand the branch endpoints, analyze the limits of
\[
r=u^{1/u}.
\]
Take logarithms:
\[
\ln r=\frac{\ln u}{u}.
\]
If \(u\to 0^+\), then \(\ln u\to-\infty\) and \((\ln u)/u\to-\infty\), so
\[
r\to 0.
\]
If \(u\to\infty\), then \((\ln u)/u\to 0\), so
\[
r\to 1.
\]
Translating back to \(\theta\),
\[
\theta\to k\pi \Longrightarrow r\to 0,
\qquad
\theta\to\frac{\pi}{2}+k\pi \Longrightarrow r\to 1.
\]
So each quadrant branch starts near the pole and approaches radius \(1\) near the vertical axis.

\item To find the largest radius, maximize \(r=u^{1/u}\). Since the logarithm is increasing, it is enough to maximize
\[
f(u)=\ln r=\frac{\ln u}{u}.
\]
Differentiate:
\[
f'(u)=\frac{1-\ln u}{u^2}.
\]
Setting \(f'(u)=0\) gives
\[
\ln u=1
\quad\Longrightarrow\quad
u=e.
\]
Therefore the maximum radius is
\[
r_{\max}=e^{1/e},
\]
and it occurs when
\[
\lvert\tan\theta\rvert=e.
\]
This tells us the bulges in each quadrant swell out to the same largest distance.

\item Because the formula depends only on \(\lvert\tan\theta\rvert\), we have
\[
r(-\theta)=r(\theta),
\qquad
r(\pi-\theta)=r(\theta).
\]
So the graph is symmetric about both coordinate axes. One complete tracing sweep is obtained by passing once through the four quadrant intervals
\[
\left(0,\frac{\pi}{2}\right),\ 
\left(\frac{\pi}{2},\pi\right),\ 
\left(\pi,\frac{3\pi}{2}\right),\ 
\left(\frac{3\pi}{2},2\pi\right).
\]
Equivalently, we can say that the whole graph is produced on
\[
0<\theta<2\pi
\]
with the axis angles omitted because the formula is undefined there.
\end{steps}
\end{solutionbox}

\begin{solutionbox}
\textbf{Method 2: Trace the curve quadrant by quadrant}

\begin{steps}
\item On the first quadrant \(0<\theta<\pi/2\), both \(\tan\theta\) and \(\cot\theta\) are positive, so
\[
r=(\tan\theta)^{\cot\theta}.
\]
As \(\theta\) increases from \(0^+\) to \(\pi/2^-\), \(\tan\theta\) runs from \(0\) to \(\infty\). This means the first-quadrant branch already contains the essential radial behavior of the whole graph.

\item Near \(\theta=0^+\), \(\tan\theta\) is tiny, so the radius tends to \(0\). Near \(\theta=\pi/2^-\), \(\tan\theta\) is huge, and the radius tends to \(1\). Therefore the first-quadrant branch starts near the pole, bulges outward, and approaches the point \((0,1)\) as a limit position.

\item In the second quadrant, \(\tan\theta<0\), but the absolute values make \(\lvert\tan\theta\rvert\) behave exactly as in the first quadrant. So the second-quadrant piece is the mirror image of the first across the \(y\)-axis. By the same reasoning, the lower half-plane is the mirror image of the upper half-plane across the \(x\)-axis.

\item Putting the quadrants together, the graph consists of four congruent bulges, one in each quadrant. The branches approach the pole along the \(x\)-axis directions and approach the points \((0,\pm1)\) as \(\theta\) approaches the vertical-axis directions. This is why the graph is commonly described as a valentine-shaped curve. A full sweep is again obtained on
\[
0<\theta<2\pi
\]
with breaks at the undefined axis angles.
\end{steps}
\end{solutionbox}

\begin{solutionbox}
\textbf{Hand-in Version}

A convenient full tracing choice is
\[
\boxed{0<\theta<2\pi}
\]
with
\[
\theta\neq \frac{\pi}{2},\ \pi,\ \frac{3\pi}{2},
\]
since the formula is undefined on those axis directions.

Let
\[
u=\lvert\tan\theta\rvert.
\]
Then
\[
r=u^{1/u}.
\]
Hence
\[
\theta\to k\pi \Longrightarrow r\to 0,
\qquad
\theta\to \frac{\pi}{2}+k\pi \Longrightarrow r\to 1,
\]
and
\[
r_{\max}=e^{1/e}
\]
when
\[
\lvert\tan\theta\rvert=e.
\]

Also,
\[
r(-\theta)=r(\theta),
\qquad
r(\pi-\theta)=r(\theta),
\]
so the graph is symmetric about both coordinate axes. Thus the curve has four matching bulges, one in each quadrant, giving the valentine shape.
\end{solutionbox}

\sectiontitle{sec:11-3}{11-3}

\problem{prob:11-3-34}{34}{Find the values of \(p\) for which the series is convergent:
\[
\displaystyle \sum_{n=1}^{\infty}\frac{\ln n}{n^p}
\]}

\begin{solutionbox}
\textbf{Method 1: Use the Integral Test}

\begin{steps}
\item Consider
\[
f(x)=\frac{\ln x}{x^p},\qquad x\ge 2.
\]
For large \(x\), this function is positive and decreasing when \(p>0\), so the Integral Test is appropriate. It converts the series question into an improper integral question:
\[
\int_2^\infty \frac{\ln x}{x^p}\,dx.
\]

\item First treat the case \(p>1\). Integrate by parts with
\[
u=\ln x,\qquad dv=x^{-p}\,dx.
\]
This choice is natural because differentiating \(\ln x\) simplifies it, while integrating \(x^{-p}\) is straightforward. Then
\[
du=\frac{dx}{x},
\qquad
v=\frac{x^{1-p}}{1-p}.
\]
So
\[
\int \frac{\ln x}{x^p}\,dx
=\frac{x^{1-p}\ln x}{1-p}-\frac{1}{1-p}\int x^{-p}\,dx.
\]
Therefore
\[
\int \frac{\ln x}{x^p}\,dx
=\frac{x^{1-p}\ln x}{1-p}-\frac{x^{1-p}}{(1-p)^2}+C.
\]
When \(p>1\), we have \(1-p<0\), so \(x^{1-p}\to 0\) as \(x\to\infty\), and also \(x^{1-p}\ln x\to 0\). Hence the improper integral converges, so the series converges.

\item Now check \(p=1\). Then the series becomes
\[
\sum_{n=1}^\infty \frac{\ln n}{n},
\]
and the corresponding integral is
\[
\int_2^\infty \frac{\ln x}{x}\,dx
=\frac{(\ln x)^2}{2}\Big|_2^\infty.
\]
This diverges because \((\ln x)^2\to\infty\). Therefore the series diverges at \(p=1\).

\item Finally, if \(p<1\), then for all sufficiently large \(n\) we have \(\ln n\ge 1\), so
\[
\frac{\ln n}{n^p}\ge \frac{1}{n^p}.
\]
The \(p\)-series \(\sum 1/n^p\) diverges when \(p\le 1\), so our series diverges by comparison. Combining all cases gives
\[
\boxed{p>1}.
\]
\end{steps}
\end{solutionbox}

\begin{solutionbox}
\textbf{Method 2: Use Cauchy Condensation}

\begin{steps}
\item For \(p\le 0\), the terms do not even go to zero:
\[
\frac{\ln n}{n^p}=(\ln n)n^{-p}.
\]
If \(p=0\), this is \(\ln n\to\infty\); if \(p<0\), it grows even faster. Since a convergent series must have terms tending to \(0\), the series diverges immediately when \(p\le 0\).

\item Now suppose \(p>0\), and let
\[
a_n=\frac{\ln n}{n^p}.
\]
For large \(n\), \(a_n\) is decreasing because the function
\[
f(x)=\frac{\ln x}{x^p}
\]
has derivative
\[
f'(x)=x^{-p-1}(1-p\ln x),
\]
which is negative once \(x>e^{1/p}\). This is exactly the condition needed to apply Cauchy Condensation.

\item The condensed series is
\[
\sum_{k=1}^\infty 2^k a_{2^k}
=\sum_{k=1}^\infty 2^k\frac{\ln(2^k)}{(2^k)^p}
=(\ln 2)\sum_{k=1}^\infty k\,2^{k(1-p)}.
\]
This step is appropriate because condensation replaces the logarithm in the numerator by the simpler factor \(k\), making the convergence threshold easier to see.

\item If \(p>1\), then \(2^{1-p}<1\), so
\[
\sum_{k=1}^\infty k\,2^{k(1-p)}
\]
is a convergent power-weighted geometric series. If \(p=1\), it becomes
\[
\sum_{k=1}^\infty k,
\]
which diverges. If \(0<p<1\), then the factor \(2^{1-p}>1\), so the terms do not even go to zero, and the series diverges. Hence condensation again gives
\[
\boxed{p>1}.
\]
\end{steps}
\end{solutionbox}

\begin{solutionbox}
\textbf{Hand-in Version}

\[
\sum_{n=1}^{\infty}\frac{\ln n}{n^p}
\]
is convergent exactly when
\[
\boxed{p>1.}
\]

Reason:
\begin{itemize}[leftmargin=*]
\item If \(p>1\), the Integral Test shows
\[
\int_2^\infty \frac{\ln x}{x^p}\,dx
\]
converges.
\item If \(p=1\), the series behaves like \(\sum (\ln n)/n\), which diverges.
\item If \(p<1\), then eventually \((\ln n)/n^p\ge 1/n^p\), and \(\sum 1/n^p\) diverges.
\end{itemize}
\end{solutionbox}

\sectiontitle{sec:11-10}{11-10}

\problem{prob:11-10-1}{1}{Use the definition of a Taylor series to find the first four nonzero terms of the series for \(f(x)\) centered at the given value of \(a\). (6 point)
\[
f(x)=\sqrt[3]{x},\ a=8
\]}

\begin{solutionbox}
\textbf{Method 1: Rewrite the function in binomial form around \(x=8\)}

\begin{steps}
\item Because the series is centered at \(a=8\), the natural variable is \(x-8\). Rewrite \(x\) as
\[
x=8+(x-8)=8\left(1+\frac{x-8}{8}\right).
\]
This is appropriate because it turns the cube root into a standard binomial expression of the form \((1+u)^{1/3}\).

\item Now compute
\[
\sqrt[3]{x}
=\left[8\left(1+\frac{x-8}{8}\right)\right]^{1/3}
=2\left(1+\frac{x-8}{8}\right)^{1/3}.
\]
Use the generalized binomial expansion
\[
(1+u)^{1/3}
=1+\frac13u+\frac{\frac13\left(-\frac23\right)}{2!}u^2
+\frac{\frac13\left(-\frac23\right)\left(-\frac53\right)}{3!}u^3+\cdots.
\]
This series is the correct tool because it is exactly the Taylor series of \((1+u)^{1/3}\) about \(u=0\).

\item Substitute
\[
u=\frac{x-8}{8}
\]
and simplify term by term:
\[
2\left(1+\frac13\frac{x-8}{8}-\frac19\frac{(x-8)^2}{8^2}
+\frac{5}{81}\frac{(x-8)^3}{8^3}+\cdots\right).
\]
Therefore
\[
\sqrt[3]{x}
=2+\frac{x-8}{12}-\frac{(x-8)^2}{288}+\frac{5(x-8)^3}{20736}+\cdots.
\]

\item These are the first four nonzero terms because the constant term and the next three powers of \(x-8\) all have nonzero coefficients.
\end{steps}
\end{solutionbox}

\begin{solutionbox}
\textbf{Method 2: Use the Taylor formula from derivatives at \(a=8\)}

\begin{steps}
\item The Taylor series centered at \(a=8\) is
\[
f(x)=f(8)+f'(8)(x-8)+\frac{f''(8)}{2!}(x-8)^2+\frac{f'''(8)}{3!}(x-8)^3+\cdots.
\]
This is appropriate because the problem explicitly asks for the Taylor series centered at \(8\).

\item Differentiate:
\[
\begin{aligned}
f(x)&=x^{1/3},\\
f'(x)&=\frac13x^{-2/3},\\
f''(x)&=-\frac29x^{-5/3},\\
f'''(x)&=\frac{10}{27}x^{-8/3}.
\end{aligned}
\]
These derivatives are easy to obtain from the power rule, so this method is straightforward.

\item Evaluate at \(x=8\). Since
\[
8^{1/3}=2,\quad 8^{2/3}=4,\quad 8^{5/3}=32,\quad 8^{8/3}=256,
\]
we get
\[
f(8)=2,
\qquad
f'(8)=\frac{1}{12},
\qquad
f''(8)=-\frac{1}{144},
\qquad
f'''(8)=\frac{5}{3456}.
\]

\item Substitute into the Taylor formula:
\[
f(x)=2+\frac{1}{12}(x-8)-\frac{1}{144}\frac{(x-8)^2}{2!}
+\frac{5}{3456}\frac{(x-8)^3}{3!}+\cdots.
\]
So
\[
\sqrt[3]{x}
=2+\frac{x-8}{12}-\frac{(x-8)^2}{288}+\frac{5(x-8)^3}{20736}+\cdots.
\]
\end{steps}
\end{solutionbox}

\begin{solutionbox}
\textbf{Hand-in Version}

The first four nonzero terms of the Taylor series of \(\sqrt[3]{x}\) about \(a=8\) are
\[
\boxed{
\sqrt[3]{x}
=2+\frac{x-8}{12}-\frac{(x-8)^2}{288}+\frac{5(x-8)^3}{20736}+\cdots
}
\]
\end{solutionbox}

\problem{prob:11-10-2}{2}{Evaluate the indefinite integral as an infinite series. (7 point)
\[
\displaystyle \int \frac{\cos x-1}{x}\,dx
\]}

\begin{solutionbox}
\textbf{Method 1: Expand the integrand first, then integrate term by term}

\begin{steps}
\item Start from the Maclaurin series for cosine:
\[
\cos x=1-\frac{x^2}{2!}+\frac{x^4}{4!}-\frac{x^6}{6!}+\cdots
=\sum_{n=0}^{\infty}\frac{(-1)^n x^{2n}}{(2n)!}.
\]
This is the right starting point because the integrand contains \(\cos x\), and power series can usually be integrated term by term inside their interval of convergence.

\item Subtract \(1\):
\[
\cos x-1
=\sum_{n=1}^{\infty}\frac{(-1)^n x^{2n}}{(2n)!}.
\]
Now divide by \(x\):
\[
\frac{\cos x-1}{x}
=\sum_{n=1}^{\infty}\frac{(-1)^n x^{2n-1}}{(2n)!}.
\]
This step is appropriate because the constant term has disappeared, so dividing by \(x\) does not create a power-series problem.

\item Integrate term by term:
\[
\int \frac{\cos x-1}{x}\,dx
=C+\sum_{n=1}^{\infty}\frac{(-1)^n}{(2n)!}\int x^{2n-1}\,dx.
\]
Thus
\[
\int \frac{\cos x-1}{x}\,dx
=C+\sum_{n=1}^{\infty}\frac{(-1)^n x^{2n}}{(2n)(2n)!}.
\]

\item Writing out the first few terms,
\[
\int \frac{\cos x-1}{x}\,dx
=C-\frac{x^2}{4}+\frac{x^4}{96}-\frac{x^6}{4320}+\frac{x^8}{322560}-\cdots.
\]
\end{steps}
\end{solutionbox}

\begin{solutionbox}
\textbf{Method 2: Assume an antiderivative power series and solve for its coefficients}

\begin{steps}
\item Let
\[
F(x)=\int \frac{\cos x-1}{x}\,dx.
\]
Assume
\[
F(x)=\sum_{n=0}^{\infty}a_n x^n.
\]
This is appropriate because the problem asks for the antiderivative \emph{as a series}, so solving directly for the coefficients is a natural alternative to integrating term by term.

\item Differentiate:
\[
F'(x)=\sum_{n=1}^{\infty} n a_n x^{n-1}.
\]
Then multiply by \(x\):
\[
xF'(x)=\sum_{n=1}^{\infty} n a_n x^n.
\]
But
\[
xF'(x)=\cos x-1
=\sum_{m=1}^{\infty}\frac{(-1)^m x^{2m}}{(2m)!}.
\]
Now both sides are power series in powers of \(x\), so coefficient comparison becomes possible.

\item Compare coefficients. There are no odd-power terms on the right, so
\[
a_{2m-1}=0 \qquad (m\ge 1).
\]
For the even powers,
\[
(2m)a_{2m}=\frac{(-1)^m}{(2m)!},
\]
so
\[
a_{2m}=\frac{(-1)^m}{(2m)(2m)!}.
\]
The constant coefficient \(a_0\) is arbitrary, so write it as \(C\).

\item Therefore
\[
F(x)=C+\sum_{m=1}^{\infty}\frac{(-1)^m x^{2m}}{(2m)(2m)!}.
\]
This matches the result from Method 1:
\[
F(x)=C-\frac{x^2}{4}+\frac{x^4}{96}-\frac{x^6}{4320}+\cdots.
\]
\end{steps}
\end{solutionbox}

\begin{solutionbox}
\textbf{Hand-in Version}

\[
\boxed{
\int \frac{\cos x-1}{x}\,dx
=C+\sum_{n=1}^{\infty}\frac{(-1)^n x^{2n}}{(2n)(2n)!}
}
\]

In expanded form,
\[
\boxed{
\int \frac{\cos x-1}{x}\,dx
=C-\frac{x^2}{4}+\frac{x^4}{96}-\frac{x^6}{4320}+\frac{x^8}{322560}-\cdots
}
\]
\end{solutionbox}

\problem{prob:11-10-93}{93}{Use the Maclaurin series for \(f(x)=x\sin(x^2)\) to find \(f^{(203)}(0)\).

\medskip
\noindent\textit{Note:} This problem uses the Maclaurin series, so the expansion is centered at \(0\). If the center were changed to some nonzero value \(a\), is there still a way to solve it by using a Taylor series centered at \(a\)?}

\begin{solutionbox}
\textbf{Method 1: Expand \(\sin(x^2)\) directly}

\begin{steps}
\item Start from the Maclaurin series
\[
\sin u=u-\frac{u^3}{3!}+\frac{u^5}{5!}-\frac{u^7}{7!}+\cdots
=\sum_{n=0}^{\infty}\frac{(-1)^n u^{2n+1}}{(2n+1)!}.
\]
This is the correct starting point because the function is \(x\sin(x^2)\), so substituting \(u=x^2\) immediately matches the structure of the problem.

\item Substitute \(u=x^2\):
\[
\sin(x^2)=x^2-\frac{x^6}{3!}+\frac{x^{10}}{5!}-\frac{x^{14}}{7!}+\cdots
=\sum_{n=0}^{\infty}\frac{(-1)^n x^{4n+2}}{(2n+1)!}.
\]
Now multiply by \(x\):
\[
f(x)=x\sin(x^2)
=x^3-\frac{x^7}{3!}+\frac{x^{11}}{5!}-\frac{x^{15}}{7!}+\cdots
=\sum_{n=0}^{\infty}\frac{(-1)^n x^{4n+3}}{(2n+1)!}.
\]

\item To find \(f^{(203)}(0)\), identify the coefficient of \(x^{203}\), because in a Maclaurin series
\[
f(x)=\sum_{m=0}^{\infty}\frac{f^{(m)}(0)}{m!}x^m.
\]
We need
\[
4n+3=203
\quad\Longrightarrow\quad
n=50.
\]
So the \(x^{203}\)-term is
\[
\frac{x^{203}}{101!}.
\]

\item Therefore
\[
\frac{f^{(203)}(0)}{203!}=\frac{1}{101!},
\]
which gives
\[
\boxed{f^{(203)}(0)=\frac{203!}{101!}}.
\]
\end{steps}
\end{solutionbox}

\begin{solutionbox}
\textbf{Method 2: Use the complex-exponential form}

\begin{steps}
\item Write
\[
\sin(x^2)=\Im\!\bigl(e^{ix^2}\bigr),
\]
so
\[
f(x)=x\sin(x^2)=\Im\!\bigl(xe^{ix^2}\bigr).
\]
This is appropriate because the exponential series is especially easy to expand and track coefficients in.

\item Expand:
\[
e^{ix^2}=\sum_{n=0}^{\infty}\frac{(ix^2)^n}{n!}
=\sum_{n=0}^{\infty}\frac{i^n x^{2n}}{n!}.
\]
Then
\[
xe^{ix^2}=\sum_{n=0}^{\infty}\frac{i^n x^{2n+1}}{n!}.
\]

\item The \(x^{203}\)-term appears when
\[
2n+1=203
\quad\Longrightarrow\quad
n=101.
\]
So the relevant term is
\[
\frac{i^{101}x^{203}}{101!}.
\]
Since
\[
i^{101}=i^{100}i=(i^4)^{25}i=i,
\]
its imaginary part is \(1\). Therefore the coefficient of \(x^{203}\) in \(x\sin(x^2)\) is again
\[
\frac{1}{101!}.
\]

\item Hence
\[
\frac{f^{(203)}(0)}{203!}=\frac{1}{101!},
\]
so
\[
\boxed{f^{(203)}(0)=\frac{203!}{101!}}.
\]
\end{steps}
\end{solutionbox}

\begin{solutionbox}
\textbf{Hand-in Version}

\[
x\sin(x^2)
=x^3-\frac{x^7}{3!}+\frac{x^{11}}{5!}-\cdots+\frac{x^{203}}{101!}+\cdots
\]

So
\[
\frac{f^{(203)}(0)}{203!}=\frac{1}{101!},
\]
hence
\[
\boxed{f^{(203)}(0)=\frac{203!}{101!}}.
\]

\medskip
\noindent\textbf{Note.} If the center were changed to some nonzero \(a\), one could still use a Taylor series centered at \(a\), but then one would have to re-expand that series into powers of \(x\) to read off the coefficient of \(x^{203}\). The Maclaurin series is the most efficient choice because it gives the coefficient at \(x=0\) directly.
\end{solutionbox}

\problem{prob:11-10-94}{94}{If \(f(x)=(1+x^3)^{30}\), what is \(f^{(58)}(0)\)?}

\begin{solutionbox}
\textbf{Method 1: Use the binomial theorem}

\begin{steps}
\item Expand:
\[
(1+x^3)^{30}
=\sum_{k=0}^{30}\binom{30}{k}(x^3)^k
=\sum_{k=0}^{30}\binom{30}{k}x^{3k}.
\]
This is the natural method because the function is already in binomial form.

\item In a Maclaurin series, the coefficient of \(x^{58}\) determines \(f^{(58)}(0)\). But every exponent in the expansion is of the form \(3k\), and
\[
58
\]
is not divisible by \(3\). Therefore there is no \(x^{58}\)-term at all.

\item Hence the coefficient of \(x^{58}\) is \(0\), so
\[
\frac{f^{(58)}(0)}{58!}=0.
\]
Therefore
\[
\boxed{f^{(58)}(0)=0}.
\]
\end{steps}
\end{solutionbox}

\begin{solutionbox}
\textbf{Method 2: Use the cube-root-of-unity symmetry}

\begin{steps}
\item Let
\[
\omega=e^{2\pi i/3}.
\]
Then \(\omega^3=1\). Because the function depends on \(x\) only through \(x^3\),
\[
f(\omega x)=(1+(\omega x)^3)^{30}=(1+\omega^3 x^3)^{30}=(1+x^3)^{30}=f(x).
\]
This symmetry is appropriate because it forces strong restrictions on which Maclaurin coefficients can be nonzero.

\item Write the Maclaurin series as
\[
f(x)=\sum_{n=0}^{\infty}a_n x^n.
\]
Then
\[
f(\omega x)=\sum_{n=0}^{\infty}a_n \omega^n x^n.
\]
Since \(f(\omega x)=f(x)\) for all \(x\), coefficient comparison gives
\[
a_n\omega^n=a_n.
\]
Thus either \(a_n=0\), or \(\omega^n=1\).

\item But \(\omega^n=1\) happens only when \(3\mid n\). So every coefficient \(a_n\) with \(n\) not divisible by \(3\) must be zero. In particular,
\[
a_{58}=0
\]
because \(58\) is not a multiple of \(3\).

\item Since
\[
a_{58}=\frac{f^{(58)}(0)}{58!},
\]
we get
\[
\boxed{f^{(58)}(0)=0}.
\]
\end{steps}
\end{solutionbox}

\begin{solutionbox}
\textbf{Hand-in Version}

\[
(1+x^3)^{30}=\sum_{k=0}^{30}\binom{30}{k}x^{3k}.
\]

All powers are multiples of \(3\), and \(58\) is not a multiple of \(3\). Therefore the coefficient of \(x^{58}\) is \(0\), so
\[
\boxed{f^{(58)}(0)=0}.
\]
\end{solutionbox}

\end{document}
